\centerline{\bf \Huge Квадроберы}

\begin{flushright}
        \scriptsize{Bombombini Gusini}
\end{flushright}

Обозначения: $\mathbb{R}, \mathbb{Q}, \mathbb{N}$ - множества действительных, рациональных, натуральных чисел соответственно (плюсик внизу добавляет условие положительности, то есть $\mathbb{R}_{+}$обозначает действительные положительные числа). Во задачах листика надо найти всевозможные функции, с заданной областью определения и областью значений, удовлетворяющих условиям задачи.

1. Попробуйте что-нибудь подставить в функциональное уравнение.

(a) $f: \mathbb{R} \rightarrow \mathbb{R}$,

$$
f(x-y)=f(x)+f(y)-2 x y
$$

(b) $f: \mathbb{R} \rightarrow \mathbb{R}$,

$$
f(x+f(y))=x+y
$$

(c) $f: \mathbb{R} \rightarrow \mathbb{R}$,

$$
x f(y)+y f(x)=(x+y) f(x) f(y)
$$

2. Надо составить систему линейных уравнений.

(a) $f: \mathbb{R} \rightarrow \mathbb{R}$,

$$
2 f(x)+f(1-x)=x^2
$$

(b) $f: \mathbb{R} \backslash\{0,1\} \rightarrow \mathbb{R}$,

$$
f(x)+f\left(\frac{1}{1-x}\right)=x
$$

3. Попробуем сделать замену переменных, функции?

(a) $f: \mathbb{R} \rightarrow \mathbb{R}_{+} \cup\{0\}$,

$$
f(x+y)=f(x)+f(y)+2 x y
$$

(b) $f: \mathbb{R} \rightarrow \mathbb{R}_{+} \cup\{0\}$,

$$
f\left(x^2+y^2\right)=f\left(x^2-y^2\right)+f(2 x y)
$$
